\section{Form HTML: Elemen form, input, label, dan button}

Formulir (form) HTML merupakan mekanisme utama untuk mengumpulkan masukan pengguna di web. Setiap kali pengguna mendaftar akun, mengirim pesan kontak, melakukan pencarian, atau mengisi survei, di balik antarmuka tersebut terdapat elemen \code{<form>}. Form memungkinkan data dari browser dikirim ke server untuk diproses, disimpan, atau ditanggapi. Pemahaman menyeluruh tentang elemen-elemen form adalah keterampilan fundamental yang dibutuhkan setiap pengembang web \cite{mdnHtmlIntro}.

Elemen \code{<form>} berfungsi sebagai wadah yang membungkus semua kontrol input. Dua atribut penting pada \code{<form>} adalah \code{action} yang menentukan URL tujuan pengiriman data, dan \code{method} yang menentukan metode HTTP yang digunakan. Metode \code{GET} mengirim data melalui URL query string (cocok untuk pencarian), sedangkan \code{POST} mengirim data di body request (lebih aman untuk data sensitif). Atribut tambahan seperti \code{enctype} diperlukan saat form mengirim file \cite{webdevHtml}.

\begin{table}[htbp]
\centering
\begin{tabular}{|P{2.5cm}|P{4cm}|P{5.5cm}|}
\hline
\textbf{Atribut} & \textbf{Nilai} & \textbf{Fungsi} \\
\hline
\code{action} & URL tujuan & Menentukan ke mana data dikirim \\
\hline
\code{method} & \code{get} / \code{post} & Metode HTTP untuk pengiriman \\
\hline
\code{enctype} & \code{multipart/form-data} & Diperlukan saat mengirim file \\
\hline
\code{novalidate} & (boolean) & Menonaktifkan validasi bawaan browser \\
\hline
\code{target} & \code{\_self}, \code{\_blank}, dll. & Di mana respons ditampilkan \\
\hline
\code{autocomplete} & \code{on} / \code{off} & Mengatur auto-fill browser \\
\hline
\end{tabular}
\caption{Atribut penting elemen \code{<form>}}
\end{table}

Elemen \code{<label>} adalah pasangan wajib untuk setiap kontrol input. Label memberikan teks deskriptif yang menjelaskan apa yang harus diisi pengguna. Pengaitan label dengan input dilakukan melalui atribut \code{for} pada label yang nilainya sama dengan \code{id} pada input. Pengaitan ini penting bukan hanya untuk aksesibilitas (pembaca layar membacakan label saat input difokuskan), tetapi juga untuk kemudahan penggunaan---klik pada label akan memfokuskan atau mengaktifkan input terkait, memperbesar area klik terutama untuk checkbox dan radio button \cite{mdnAccessibility}.

\begin{figure}[htbp]
\centering
\resizebox{\textwidth}{!}{%
\begin{tikzpicture}[node distance=2cm, transform shape, every node/.style={font=\footnotesize}]
  \node (user) [class, fill=orange!20] {Pengguna\\mengisi form};
  \node (browser) [class, right=of user] {Browser\\validasi client};
  \node (send) [class, right=of browser] {HTTP Request\\(GET/POST)};
  \node (server) [class, right=of send] {Server\\memproses data};
  \node (response) [class, right=of server] {Response\\(HTML/redirect)};

  \draw [arrow] (user) -- (browser);
  \draw [arrow] (browser) -- node[above, font=\scriptsize] {valid} (send);
  \draw [arrow, dashed] (browser) to[bend right=30] node[below, font=\scriptsize] {invalid: tampil error} (user);
  \draw [arrow] (send) -- (server);
  \draw [arrow] (server) -- (response);
  \draw [arrow, dashed] (response) to[bend right=25] node[below, font=\scriptsize] {halaman baru} (user);
\end{tikzpicture}%
}
\caption{Alur pengiriman data form: dari pengguna ke server dan kembali}
\end{figure}

Elemen \code{<button>} digunakan untuk membuat tombol interaktif di dalam form. Terdapat tiga tipe button: \code{type="submit"} yang mengirim data form ke server (perilaku default jika \code{type} tidak ditentukan), \code{type="reset"} yang mengosongkan semua isian form ke nilai awal, dan \code{type="button"} yang tidak melakukan aksi form tetapi dapat diprogram dengan JavaScript. Sebaiknya selalu tentukan atribut \code{type} secara eksplisit untuk menghindari perilaku yang tidak diinginkan \cite{whatwgHtml}.

\begin{webcode}[language=HTML, caption={Form kontak lengkap dengan label dan button}]
<form action="/kirim-pesan" method="post">
  <p>
    <label for="nama">Nama Lengkap:</label>
    <input type="text" id="nama" name="nama"
           placeholder="Masukkan nama" required />
  </p>
  <p>
    <label for="email">Alamat Email:</label>
    <input type="email" id="email" name="email"
           placeholder="nama@domain.com" required />
  </p>
  <p>
    <label for="pesan">Pesan:</label>
    <textarea id="pesan" name="pesan" rows="5"
              placeholder="Tulis pesan Anda..."></textarea>
  </p>
  <p>
    <button type="submit">Kirim Pesan</button>
    <button type="reset">Batal</button>
  </p>
</form>
\end{webcode}

\begin{konsep}
\textbf{Selalu Gunakan Label untuk Setiap Input.} Form tanpa label yang tepat menyulitkan pengguna pembaca layar dan mengurangi kemudahan penggunaan. Dua cara mengaitkan label: (1) \textbf{Eksplisit}: \code{<label for="email">} dengan \code{<input id="email">}; (2) \textbf{Implisit}: membungkus input di dalam label, misalnya \code{<label>Email: <input type="email"></label>}. Cara eksplisit lebih direkomendasikan karena lebih fleksibel dalam penataan layout \cite{webdevAccessibility}.
\end{konsep}

